% RQ2 Theoretical Framework: Governance Capture

We develop formal metrics for governance capture and analyze the theoretical properties of mitigation mechanisms.

\subsection{Capture Metrics}

\subsubsection{Whale Influence Index}

We define whale influence as the fraction of voting power held by the top $k$ addresses:

\begin{equation}
\text{WhaleInfluence}_k = \frac{\sum_{i=1}^{k} w(a_{(i)})}{\sum_{j=1}^{n} w(a_j)}
\end{equation}

where $a_{(i)}$ is the $i$-th largest holder. We typically use $k=5$ or $k=10$.

\subsubsection{Nakamoto Coefficient}

The minimum number of actors needed to control majority voting power:

\begin{equation}
\text{Nakamoto} = \min \left\{ k : \sum_{i=1}^{k} w(a_{(i)}) > 0.5 \cdot \sum_{j=1}^{n} w(a_j) \right\}
\end{equation}

Higher values indicate more distributed power.

\subsubsection{Gini Coefficient of Voting Power}

Standard inequality measure applied to voting weights:

\begin{equation}
\text{Gini} = \frac{\sum_{i=1}^{n} \sum_{j=1}^{n} |w(a_i) - w(a_j)|}{2n \sum_{i=1}^{n} w(a_i)}
\end{equation}

\subsubsection{Capture Risk Score}

We define a composite capture risk:

\begin{equation}
\text{CaptureRisk} = \frac{\text{WhaleInfluence}_{10}}{\text{Nakamoto}} \cdot (1 + \text{Gini})
\end{equation}

This captures both concentration (whale influence) and fragility (low Nakamoto coefficient).

\subsection{Mitigation Mechanism Analysis}

\subsubsection{Vote Caps}

With cap $c$, the effective voting power is:

\begin{equation}
w_c(a) = \min(T(a), c)
\end{equation}

Properties:
\begin{itemize}
    \item Maximum individual influence bounded by $c$
    \item Creates ``cliff effect'' at threshold
    \item Sybil attack cost: splitting $n$ tokens across $\lceil n/c \rceil$ addresses
    \item Gas cost of sybil: $O(n/c)$ addresses to manage
\end{itemize}

\subsubsection{Quadratic Voting}

Voting power scales as square root:

\begin{equation}
w_q(a) = \sqrt{T(a)}
\end{equation}

Properties:
\begin{itemize}
    \item No hard cutoffs; smooth diminishing returns
    \item Whale with $100\times$ tokens has only $10\times$ power
    \item Sybil resistant to first order: $\sqrt{a} + \sqrt{b} < \sqrt{a+b}$ only when $ab < 0$
    \item However: multiple identities still beneficial under quadratic
\end{itemize}

\subsubsection{Velocity Penalties}

Tokens acquired within window $\tau$ receive reduced weight:

\begin{equation}
w_v(a, t) = T_{\text{old}}(a, t) + \alpha \cdot T_{\text{new}}(a, t, \tau)
\end{equation}

where:
\begin{itemize}
    \item $T_{\text{old}}(a, t) = $ tokens held before $t - \tau$
    \item $T_{\text{new}}(a, t, \tau) = $ tokens acquired in $[t-\tau, t]$
    \item $\alpha \in [0, 1)$ is the penalty factor (0 = full penalty, 1 = no penalty)
\end{itemize}

Properties:
\begin{itemize}
    \item Prevents ``flash governance'' attacks
    \item Rewards long-term holders
    \item Does not address existing concentration
    \item Complexity: requires tracking acquisition times
\end{itemize}

\subsection{Theoretical Predictions}

\subsubsection{Prediction 1: Cap Effectiveness}

Vote caps reduce whale influence proportionally until sybil attacks become economically attractive:

\begin{equation}
\text{WhaleInfluence}(c) \approx \min\left(\text{WhaleInfluence}_0, \frac{k \cdot c}{\text{TotalSupply}}\right)
\end{equation}

\subsubsection{Prediction 2: Quadratic Compression}

Quadratic voting compresses the Gini coefficient:

\begin{equation}
\text{Gini}_{\text{quad}} \approx 0.5 \cdot \text{Gini}_{\text{token}}
\end{equation}

\subsubsection{Prediction 3: Velocity Window Optimum}

There exists an optimal velocity window $\tau^*$ balancing attack prevention and legitimate participation:
\begin{itemize}
    \item $\tau$ too short: attacks succeed with minimal delay
    \item $\tau$ too long: legitimate new participants disadvantaged
\end{itemize}

\subsection{Throughput Tradeoff}

Mitigation mechanisms may reduce governance throughput if:
\begin{itemize}
    \item Caps reduce total voting power below quorum requirements
    \item Quadratic voting increases effective quorum difficulty
    \item Velocity penalties exclude recently-onboarded active participants
\end{itemize}

We measure this via proposal pass rate and time-to-decision.
